\begin{figure}[htbp]
\centering
\begin{tikzpicture}[causal,x=1.25cm,y=1.1cm]
 \draw (-3.6,-1.65) rectangle (-0.1,1.65);
 \node[align=left,font=\small] at (-1.85,.45) {$X_1:=f_1(X_3,N_1)$\\$X_2:=f_2(X_1,N_2)$\\$X_3:=f_3(N_3)$\\$X_4:=f_4(X_2,X_3,N_4)$};
 \node[align=left,font=\small] at (-1.85,-1.15) {$\bullet$ $N_1,\ldots,N_4$ 相互独立\\$\bullet$ $\mathcal G$ 无环};
 \draw (0,-1.65) rectangle (3.5,1.65);
 \node at (.7,1.15) {$\mathcal G$};
 \node[vtx] (ch06a-x1) at (1.75,1.1) {$X_1$};
 \node[vtx] (ch06a-x2) at (.8,0) {$X_2$};
 \node[vtx] (ch06a-x3) at (2.7,0) {$X_3$};
 \node[vtx] (ch06a-x4) at (1.75,-1.1) {$X_4$};
 \draw[ed] (ch06a-x3)--(ch06a-x1);
 \draw[ed] (ch06a-x1)--(ch06a-x2);
 \draw[ed] (ch06a-x2)--(ch06a-x4);
 \draw[ed] (ch06a-x3)--(ch06a-x4);
\end{tikzpicture}
\caption{结构因果模型的一个例子（左）及其对应图（右）。因果排序 $\pi$ 只有一种，即满足 $3\mapsto1$、$1\mapsto2$、$2\mapsto3$、$4\mapsto4$ 的排序。}
\end{figure}
